\begin{longtable}{|p{3.5cm}|p{5.5cm}|p{6cm}|}

\hline

Categoria & Características dos usuários e contextos identificados na literatura & Exemplos de estudos \\ \hline

\endfirsthead

\hline

Categoria & Características dos usuários e contextos identificados na literatura & Exemplos de estudos \\ \hline

\endhead

\hline

\endfoot

\endlastfoot

Faixa etária e perfil dos usuários &

Idade dos participantes, principalmente pessoas idosas, adultos mais velhos e usuários novatos, além de comparações entre diferentes faixas etárias e entre usuários idosos e jovens. &

Age related differences and the depth vs. breadth tradeoff in hierarchical online information systems; Redesigning web sites for older adults; Evaluating websites for older adults: adherence to 'senior-friendly' guidelines and end-user performance. \\ \hline

Experiência e familiaridade com tecnologia &

Experiência com computadores, experiência com a Web, frequência de uso, familiaridade com tecnologias, nível de conhecimento de Internet, experiência prévia com sistemas e distinção entre usuários novatos e experientes. &

Exploring the wiki user experience: The effects of training spaces on novice user usability and anxiety towards wiki editing; Investigating the roles of knowledge and cognitive abilities in older adult information seeking on the Web; Access to Legal Services: Organizing Better Self-help Systems. \\ \hline

Alfabetização informática e proficiência digital &

Baixa alfabetização informática, baixa alfabetização digital, pouca experiência com computadores, ausência de treinamento formal, diferentes níveis de proficiência tecnológica e necessidade de apoio durante a aprendizagem do uso do sistema. &

Adding voices to support web navigation among a low digital literacy group; Exploring the wiki user experience: The effects of training spaces on novice user usability and anxiety towards wiki editing; A web accessing tool for blind and visually impaired people using Bahasa Indonesia. \\ \hline

Características cognitivas &

Declínio cognitivo, memória, atenção, aprendizagem, conhecimento estrutural, capacidade de compreensão, modelos mentais, ansiedade e dificuldades relacionadas à complexidade das tarefas. &

Investigating the roles of knowledge and cognitive abilities in older adult information seeking on the Web; Health problem solving by older persons using a complex government web site: Analysis and implications for web design; Addressing Age-Related Accessibility Needs of Senior Users Through Model-Driven Engineering. \\ \hline

Características visuais &

Declínio visual, baixa acuidade visual, presbiopia, daltonismo, necessidade de fontes maiores, dificuldades de leitura e necessidades relacionadas ao contraste e à apresentação das informações. &

Personalized service studies CSS for web accessibility improvement of the senior generation; Interface personalization through inclusive user modelling web service; Research on Web Interface barrier-free Design for Elderly People. \\ \hline

Características motoras e físicas &

Limitações motoras, tremores, força de preensão, amplitude de movimento, dificuldades de interação física e outras limitações funcionais relacionadas à idade ou a deficiências. &

Breaking Barriers: The Design and Development of an Assistive Technology Web App for Older Latinos with Disabilities in Daily Activities; Interface personalization through inclusive user modelling web service; Using an Error Detection Strategy for Improving Web Accessibility for Older Adults. \\ \hline

Características auditivas e de comunicação &

Deficiências auditivas, necessidades relacionadas à comunicação, uso de voz, leitura de informações em voz alta e necessidade de recursos alternativos para apresentação do conteúdo. &

A web accessing tool for blind and visually impaired people using Bahasa Indonesia; Designing Inclusive Smartwatch Interfaces: Guidelines for Enhancing Usability and Adoption Among Older Adults; Addressing Age-Related Accessibility Needs of Senior Users Through Model-Driven Engineering. \\ \hline

Deficiências e necessidades de acessibilidade &

Deficiências visuais, auditivas, motoras e cognitivas, necessidades de tecnologias assistivas, entrada alternativa, leitura em voz alta, reorganização do conteúdo e adaptação da interface às necessidades individuais. &

A methodological approach for designing and developing web-based inventories of mobile Assistive Technology applications; Breaking Barriers: The Design and Development of an Assistive Technology Web App for Older Latinos with Disabilities in Daily Activities; A Framework for Serving Inclusive Web Forms to Disabled and Elderly Individuals. \\ \hline

Características sociodemográficas e culturais &

Gênero, escolaridade, ocupação, renda, condições médicas, contexto cultural, experiência educacional e outras características sociodemográficas consideradas nas avaliações. &

Old and Young Users' White Space Preferences for Online News Web Pages; Breaking Barriers: The Design and Development of an Assistive Technology Web App for Older Latinos with Disabilities in Daily Activities; Applicability of Process Mining in Usability Tests: A Case Study for Identifying User Mental Models in Geospatial Search Engines. \\ \hline

Contexto e domínio de uso &

Contextos como saúde, serviços governamentais, serviços tributários, bancos, turismo, educação, atividades cotidianas, busca de informações, redes sociais e uso doméstico de tecnologias. &

Health problem solving by older persons using a complex government web site: Analysis and implications for web design; Family-centered design: Interactive performance testing and user interface evaluation of the Slovenian eDavki public tax portal; Seniors' experiences with online banking. \\ \hline

Dispositivo e ambiente de interação &

Computadores, dispositivos móveis, smartwatches, interfaces Web, ambientes domésticos, dispositivos antigos e diferentes condições de acesso e utilização das tecnologias. &

Effectiveness of Using Web Applications to Preserve Cognitive Functionality in Older Adults: Mobile First Experience; Designing Inclusive Smartwatch Interfaces: Guidelines for Enhancing Usability and Adoption Among Older Adults; User Centered Design Approach for Elderly People in using Website. \\ \hline

\caption{Características dos usuários e contextos de uso identificados nos estudos que respondem diretamente à RQ4.}

\label{tab:caracteristicasRQ4}

\end{longtable}